\documentclass[12pt,a4paper]{report}
\usepackage[utf8]{inputenc}
\usepackage[T1]{fontenc}
\usepackage[french]{babel}
\usepackage{geometry}
\usepackage{graphicx}
\usepackage{lmodern}
\usepackage{setspace}
\usepackage{fancyhdr}
\pagestyle{fancy}

% Corriger ce que contient \leftmark pour les chapitres
\renewcommand{\chaptermark}[1]{%
  \markboth{CHAPITRE \thechapter. \ #1}{}%
}

\fancyhf{} % Réinitialise tout
\fancyhead[L]{\nouppercase{\leftmark}} % Titre du chapitre
\fancyfoot[L]{Evan COMBOT}
\fancyfoot[C]{\thepage}
\fancyfoot[R]{Université de Montpellier}

% Appliquer aussi aux pages de style "plain"
\fancypagestyle{plain}{%
  \fancyhf{}
  \fancyhead[L]{\nouppercase{\leftmark}}
  \fancyfoot[L]{Evan COMBOT}
  \fancyfoot[C]{\thepage}
  \fancyfoot[R]{Université de Montpellier}
}

\usepackage{titling}
\usepackage{mathpazo}
\usepackage{ragged2e}
\usepackage{titlesec}
\titleformat{\chapter}[display]{\normalfont\huge\bfseries}{\chaptername\ \thechapter}{10pt}{\Huge}
\titlespacing*{\chapter}{0pt}{-20pt}{30pt}
\setcounter{tocdepth}{2}

% === Pour gérer les marges58 ===
\geometry{left=2.0cm, right=2.0cm, top=2.0cm, bottom=2.0cm}
\onehalfspacing
\usepackage[]{hyperref}


\begin{document}
\justifying

\begin{titlepage}
\begin{center}

% === LOGOS ALIGNÉS EN LIGNE AVEC LÉGENDES EN GRAS ===
\vspace*{0.5cm}
\begin{minipage}{0.3\textwidth}
    \centering
    \includegraphics[width=0.85\textwidth]{Assets/UM.png}\\[0.2cm]
    {\footnotesize\textbf{Université de Montpellier}}
\end{minipage}
\hfill
\begin{minipage}{0.3\textwidth}
    \centering
    \includegraphics[width=0.85\textwidth]{Assets/Dpt_Info.png}\\[0.2cm]
    {\footnotesize\textbf{Département Informatique -- IMAGINE}}
\end{minipage}
\hfill
\begin{minipage}{0.3\textwidth}
    \centering
    \includegraphics[width=0.85\textwidth]{Assets/CEA.png}\\[0.2cm]
    {\footnotesize\textbf{CEA -- Commissariat à l'énergie atomique et aux énergies alternatives}}
\end{minipage}

\vspace{1.8cm}

% === ÉTABLISSEMENT & DIPLÔME ===
{\Large \textbf{Université de Montpellier}}\\[0.2cm]
{\normalsize Faculté des Sciences}\\[0.3cm]
{\large \textbf{Master 2 Informatique -- Parcours \og IMAGINE \fg{}}}\\[1.5cm]

% === TITRE DU DOCUMENT ===
{\Huge \bfseries Rapport de Stage}\\[0.4cm]
\rule{0.65\linewidth}{0.7pt}\\[1cm]

{\large \textbf{Exploitation d'images de caractérisation de chaîne radiographique}}\\[0.5cm]
\textit{Stage réalisé au sein du \textbf{Commissariat à l'énergie atomique et aux énergies alternatives}}\\[0.3cm]
\textit{du \textbf{3 février 2025} au \textbf{1 août 2025}}\\[2cm]

% === INFOS ÉTUDIANT & ENTREPRISE ===
\begin{flushleft}
\textbf{Réalisé par :} \hfill COMBOT Evan\\[0.15cm]
\textbf{Année universitaire :} \hfill 2024--2025\\[0.15cm]
\textbf{Tuteur en entreprise :} \hfill Nom Prénom\\[0.15cm]
\textbf{Encadrant pédagogique :} \hfill PUECH William\\[0.15cm]
\textbf{Adresse de l'entreprise :} \hfill CEA Gramat, 46500 Gramat, France
\end{flushleft}

\vfill
{\small Date de rendu : 20 juin 2025}

\end{center}
\end{titlepage}

%======================
% Remerciements
%======================
\chapter*{Remerciements}
\addcontentsline{toc}{chapter}{Remerciements}

Je souhaite exprimer ma profonde gratitude à toutes les personnes qui m'ont soutenu(e) et accompagné(e) tout au long de ce stage.

Je remercie tout particulièrement :
\begin{itemize}
  \item M. Nom Prénom, mon tuteur de stage en entreprise, pour son accueil chaleureux, son encadrement attentif et sa grande disponibilité ;
  \item M. William Puech, mon encadrant pédagogique, pour son suivi rigoureux et sa bienveillance;
  \item L’ensemble de l’équipe du CEA, pour leur soutien constant, leur convivialité, l’excellente ambiance de travail qu’ils ont su instaurer, ainsi que pour m’avoir intégré(e) très rapidement au sein de leur équipe.
\end{itemize}

Je tiens également à remercier l’ensemble des enseignants du Master 2 IMAGINE pour la qualité de l’enseignement et les compétences qu’ils m’ont permis d’acquérir.

\vspace{1em}
Merci à toutes et à tous.


\tableofcontents
\pagestyle{fancy}
\newpage


\chapter{Introduction}
\section{Contexte du stage}
\section{Objectifs du rapport}
Ce rapport à pour objectif de présenter le déroulement de mon stage, réalisé au sein du CEA dans le cadre du stage de fin d'études du MAster 2 IMAGINE.

Plus précisément, ce document a pour buts de :
\begin{itemize}
  \item Expliquer le cadre général du stage et les objectifs fixés en début de mission ;
  \item Décrire les différentes tâches réalisées et les méthodologies employées ;
  \item Mettre en évidence les outils et technologies utilisés dans le contexte professionnel ;
  \item Présenter les résultats obtenus et les éventuelles difficultés rencontrées ;
  \item Faire un bilan personnel et professionnel sur les acquis du stage.
\end{itemize}

L'ensemble de ce rapport permet également la mise en pratique des connaissances théoriques acquises tout au long de la formation tout en intégrant les acquis du stage dans le développement de nouvelles compétences.

... % contenu identique

\chapter{Présentation de l'entreprise}
\section{Le Commissariat à l'énergie atomique et aux énergies alternatives}
\subsection{Historique}

L'origine du Commissariat à l'énergie atomique (CEA) remonte à novembre 1895, lorsque le physicien allemand Wilhelm Conrad Röntgen découvre un rayonnement invisible et inconnu qu'il appelle \textit{rayon X}. En 1896, Henri Becquerel, un physicien français, découvre le phénomène de la radioactivité naturelle en observant une émission de rayons, analogues aux rayons X, dans les composés de l'uranium. Deux ans plus tard, Pierre et Marie Curie démontrent que certains minerais naturels émettent plus de rayons que d'autres et mettent en évidence deux éléments très radioactifs : le polonium et le radium. En 1934, Irène et Frédéric Joliot-Curie « imitent » la nature en créant un nouveau corps radioactif, le phosphore-30. Ils découvrent ainsi la radioactivité artificielle et pensent qu'en cassant des atomes, il serait possible de provoquer des réactions en chaîne explosives, libérant une énorme quantité d’énergie. C'est ce qu'on appelle la \textit{fission nucléaire}. Après la Seconde Guerre mondiale, le général de Gaulle est convaincu de la nécessité de fonder un organisme national consacré à l'énergie nucléaire. Par l'ordonnance du 18 octobre 1945, le \textit{Commissariat à l'énergie atomique (CEA)} est créé. En 2009, le CEA étend son champ d'action pour inclure les énergies renouvelables et devient officiellement le \textit{Commissariat à l'énergie atomique et aux énergies renouvelables (CEA)}.

\subsection{Activités}
Actuellement, le \textit{Commissariat à l'énergie atomique et aux énergies alternatives (CEA)} est structuré autour de quatre grandes directions opérationnelles, chacune ayant un champ de compétences spécifique :

\begin{itemize}
  \item \textbf{La Direction des énergies (DES)} : chargée des recherches dans le domaine de l’énergie nucléaire civile, des systèmes énergétiques innovants et de la transition énergétique.
  
  \item \textbf{La Direction des applications militaires (DAM)} : responsable des missions liées à la défense nationale, notamment la dissuasion nucléaire et la sécurité des installations.

  \item \textbf{La Direction de la recherche technologique (DRT)} : orientée vers l’innovation technologique et les partenariats industriels, en particulier dans les domaines de la microélectronique, des systèmes intelligents et des énergies renouvelables.

  \item \textbf{La Direction de la recherche fondamentale (DRF)} : dédiée à l’exploration des sciences fondamentales, telles que la physique nucléaire, la biologie, la climatologie ou encore les sciences de l’univers.
\end{itemize}

\subsection{Organisation}
Aujourd'hui, le CEA est composé de neuf centres de recherches dont six sont militaires, dirigé par la DAM et cinq sont destinés pour des applications civiles. C'est également cinquante-deux unités mixtes de recherche (UMR), 725 dépôts de brevets prioritaires, vingt-six Équipex (équipements d'excellence), de seize Labex (laboratoires d'excellence) et un effectif total de 16564 employés en 2023.


\section{La Direction des Applications Militaires}
\subsection{Historique \& activités}
La direction des applications militaires est créée secrètement en 1955 sous le nom de Bureau d'étude générales (BEG). Ce n'est qu'en 1958 que le génréral De Gaulle décide de créer la Direction des applications militaires, la DAM, au sein du Commissariat à l’Energie Atomique. Les armes nucléaires, les réacteurs nucléaires pour la propulsion navale et les matières nucléaires stratégiques sont au cœur de sa mission pour la dissuasion nucléaire française. Aujourd'hui, la DAM permet de contribuer à l'excellence de la recherche, à la compétitivité de l'industrie, à surveiller, analyser et intervenir pour la défense et la sécurité et également de répondre aux enjeux de la dissuasion nucléaire. Aujourd'hui la Direction des applications militaires comptait près de 4 802 employés en 2023.
\subsection{Le centre de Bruyères-le-Châtel}



\chapter{Présentation de la mission}
\section{Contexte de la mission}
Contexte professionnel de ta mission, ce que tu devais faire, les enjeux, la méthodologie, les objectifs pratiques.

qui s'incrit dans le programme SImulation
\section{Enjeux et objectifs}
... % à compléter

\chapter{Environnement technique}
\section{Technologies utilisées}
\section{Outils de développement}
\section{Contraintes techniques}
... % à compléter

\chapter{Veille technologique}
Ce que tu as appris ou approfondi pour comprendre les bases scientifiques ou techniques liées à l’outil/la machine/le principe physique utilisé dans ta mission.

Explications machine + phsique théorique sur les onde
Imagerie par Rayons X et par neutrons

... % à compléter

\chapter{Missions réalisées}
\section{Tâches effectuées}
\section{Méthodologie}
\section{Difficultés rencontrées}
... % contenu identique avec tâches, technos, outils, etc.


\section{Compétences transversales}
Par ailleurs, à l’occasion d’une formation interne organisée en début de stage, j’ai eu l’opportunité de visiter ... centre calcul

\subsection{Compétences professionnelles principales}
\subsection{Compétences acquises par des formations}


... % contenu identique

\chapter{Conclusion}
\section{Bilan du stage}
\section{Perspectives}
... % contenu identique

\appendix
\chapter{Annexes}


\end{document}
